% !TeX spellcheck = he_IL
%% ! TEX TS-program = xelatexmk
% above: commands appropriate for TeXShop in mac OS
\documentclass[12pt, a4paper]{article}
\usepackage{etoolbox}
\usepackage{bbm}
%\usepackage{pifont}
\usepackage{tikz}
%\input{preamble/UnitPreamble}
%%%%%%%%%%%%%%%%%%%%%%%%%%
%%
%%
%%  Exam Template
%%
%%
%%%%%%%%%%%%%%%%%%%%%%%%%%
\newtoggle{testreview}
\toggletrue{testreview}  
\togglefalse{testreview} % uncomment to make it ready for exam unit
\newtoggle{testsolution}
\toggletrue{testsolution}  
\togglefalse{testsolution} % comment to show solutions
\newcounter{version}
\setcounter{version}{2}  % version of multiple-choice

	





\usepackage{graphicx}

\usepackage{braket}
\usepackage[a4paper, top=3cm,left=2.5cm, right=2.3cm, bottom=2.7cm]{geometry}
\usepackage{amsmath, amssymb, array, multicol, enumitem}
\usepackage{pdfpages}
\usepackage{comment, graphicx}  
\usepackage{xparse}
\DeclareGraphicsExtensions{.png,.jpg}
\usepackage{pgffor}
\usepackage[colorlinks=true]{hyperref} 
\AtBeginDocument{%
	\hypersetup{pdfborder={0 0 0},
		urlcolor=blue}%
}

\usepackage{wasysym}
\usepackage{fontspec}
\usepackage{xunicode}
\usepackage{polyglossia}
\setdefaultlanguage{hebrew}
\setotherlanguage{english}
%\newfontfamily\hebrewfont[Script=Hebrew, Ligatures=TeX, Scale=1.1, WordSpace =1.05]{David}
%\newfontfamily\hebrewfont[Script=Hebrew, Ligatures=TeX, Scale=1.1, WordSpace=1.05]{David  CLM}
\newfontfamily\hebrewfont[Script=Hebrew]{Times New Roman}
%\newfontfamily\hebrewfont[Script=Hebrew]{Ezra SIL}

% Hebrew monospace font (must support Hebrew)
\newfontfamily\hebrewfonttt[Script=Hebrew]{FreeMono}

%http://tex.stackexchange.com/questions/141231/equation-number-enclosing-parentheses-misplaced-with-polyglossiahebrewamsmath
\makeatletter
%\def\maketag@@@#1{\hbox{\m@th\normalfont\LRE{#1}}}
\def\tagform@#1{\maketag@@@{(\ignorespaces#1\unskip)}}
\makeatother

%
% Hebrew fonts are two small 
%
\fontsize{14}{15}

%
% Package fro proper positioning of logo
%
\usepackage[absolute]{textpos}


% http://tex.stackexchange.com/questions/10589/fix-nested-section-numbers-in-rtl-languages-with-polyglossia/
\makeatletter
\renewcommand{\SepMark}[1]{\def\@SepMark{#1}}
\makeatother
\SepMark{.}
\robustify\hebrewnumeral
\robustify\Hebrewnumeral
\robustify\Hebrewnumeralfinal


% from polyglossia docs
%\gappto\captionshebrew{%
\def\urladdrname{{\itshape כתובת רשת}}
\def\emailaddrname{{\itshape כתובת דואל}}
%}

% https://tex.stackexchange.com/questions/176661/polyglossia-arabic-text-english-bibliography
\makeatletter
\let\orig@MakeUppercase \MakeUppercase
\makeatother
\providecommand{\MakeUppercase}[1]{#1}




%%%%%%%%%%%%%%%%%%%%%%%%%%
%%
%%
%%   User Commands
%%
%%
%%%%%%%%%%%%%%%%%%%%%%%%%%
\newcommand{\heb}[1]{\begin{small}\R{#1}\end{small}}
\newcommand{\hebt}[1]{\begin{scriptsize}\R{#1}\end{scriptsize}}
\newcommand{\cc}{\mathbb{C}}
\newcommand{\nn}{\mathbb{N}}
\newcommand{\zz}{\mathbb{Z}}
%\renewcommand{\phi}{\varphi}
\newcommand{\abs}[1]{\left|#1\right|}
\newcommand{\lr}[1]{\left(#1\right)}


\newcommand{\ff}{\mathbb{F}}
\newcommand{\FF}{\mathbb{F}}
\newcommand{\mtn}{m\times n}
\newcommand{\longr}{\longrightarrow}
\newcommand{\longrl}{\longleftrightarrow}
\newcommand{\ol}{\overline}
\newcommand{\ntn}{n\times n}
\newcommand{\rr}{\mathbb{R}}
\newcommand{\QQ}{\mathbb{Q}}
\newcommand{\RR}{\mathbb{R}}
\newcommand{\im}{\mathrm{Im}}

\newcommand{\CC}{\mathbb{C}}
\newcommand{\ZZ}{\mathbb{Z}}
\newcommand{\qq}{\mathbb{Q}}
\DeclareMathOperator{\vol}{vol}

\DeclareMathOperator{\Adj}{Adj}
\DeclareMathOperator{\deter}{det}
\DeclareMathOperator{\Img}{Im}
\DeclareMathOperator{\Ker}{Ker}
\DeclareMathOperator{\rank}{rank}
\DeclareMathOperator{\Span}{Span}
\DeclareMathOperator{\tr}{tr}


\newcommand{\cB}{\mathcal{B}}
\newcommand{\cC}{\mathcal{C}}
\newcommand{\cR}{\mathcal{R}}
\newcommand{\cD}{\mathcal{D}}
\newcommand{\calE}{\mathcal{E}}


%%%% new for complex functions

\newcommand{\Arg}{\mathrm{Arg}}
\newcommand{\wt}{\widetilde}
\newcommand{\wh}{\widehat}
\newcommand{\re}{\mathrm{Re}}
\newcommand{\goes}{\longrightarrow}
\newcommand{\cO}{\mathcal{O}}
\newcommand{\Log}{\mathrm{Log}}


\newcommand{\twopartdef}[4]
{
	\left\{
		\begin{array}{ll}
			#1 & \mbox{   } #2 \\
			#3 & \mbox{   } #4
		\end{array}
	\right.
}

\newcommand{\solutionword}{\textbf{פתרון:}}
\newcommand{\points}[1]{(#1 נק') }
\iftoggle{testsolution}{%
	\newcommand{\sols}[1]{\solutionword$\;\;$ #1$\square$}
	\newcommand{\spacing}{0cm}
	\newcommand{\questionspacing}{2.5cm}
	\newcommand{\extrapages}{}
}{ %
%\newcommand{\spacing}{12cm}
%\newcommand{\halfspacing}{0.5cm}	
\newcommand{\spacing}{24cm}
\newcommand{\questionspacing}{\spacing}
\newcommand{\sols}[1]{}
\newcommand{\extrapages}{%
	\newpage
	\noindent
	דף זה ריק במכוון ומיועד לשימוש הסטודנט/ית.
	\newpage
	\noindent
	דף זה ריק במכוון ומיועד לשימוש הסטודנט/ית.
}
\iftoggle{testreview}{%
\renewcommand{\spacing}{1cm}
%\newcommand{\questionspacing}{\spacing}
\renewcommand{\extrapages}{}	
}{}	
}


\DeclareMathOperator{\row}{Row}

\DeclareMathOperator{\col}{Col}


%%%%%%%%%%%%%%%%%%%%%%%%
%%
%%   BEGIN DOCUMENT
%%
%%%%%%%%%%%%%%%%%%%%%%%%

\begin{document}

\selectlanguage{hebrew}


\newcommand{\semester}{%
	סמסטר א' תשפ״ו
}
\newcommand{\course}{%
 אלגברה לינארית להנדסה 
}
\newcommand{\coursenumber}{9321}	
\newcommand{\lecturers}{%
	מרצים: ד״ר י.מייזל, ד״ר נ.גולקו, ד״ר נ.מאיר, ד״ר נ.איידלשטין, ד״ר מ.פורת, ד״ר מ.זיו-אב, ד״ר מ.ברשדסקי
}

\newcommand{\testdate}{{23.2.26}}
\iftoggle{testsolution}{%
	\newcommand{\testtype}{%
		פתרונות למבחן מועד א'
	}
}{%
	\newcommand{\testtype}{%
	\boxed{\text{
	מועד ב'
	}	}
}}

\newcommand{\duration}{%
	3 שעות.
}
\newcommand{\helpmaterial}{%
מחשבון פשוט בלבד.
}



\thispagestyle{empty}

\begin{textblock*}{9cm}[0.5,0.5](90mm,3cm)
	\includegraphics[width=5cm]{logo.jpg}
\end{textblock*}	

\noindent

\begin{center}
אוניברסיטת בן-גוריון בנגב - היחידה להוראת קורסי יסוד - %
\smallskip
\semester \\
{\large \bf \course - 214.1.\coursenumber}\\
\smallskip
\lecturers\\
\smallskip
\smallskip
\smallskip
{\Large \testtype }\\
\end{center}

\noindent
\textbf{%
תאריך:
} \testdate \quad 
\textbf{משך הבחינה: 
} \duration \quad
\textbf{
חומר עזר: 
}
\helpmaterial \\
דף הנוסחאות מצורף למבחן. אין להכניס כל חומר עזר נוסף.
\subsection*{חלק א': שאלות רב ברירה}
\begin{itemize}
    \item בחלק א' של המבחן יש 8 שאלות רב-ברירה, ומשקל כל שאלה 5 נקודות.
    \item לכל שאלה יש תשובה אחת נכונה בדיוק. יש לסמן תשובה אחת לכל שאלה \textbf{בטבלה בדף זה}, אין צורך בנימוק.
    \item שימו לב להנחיות מילוי הטבלה בעמוד הבא.
\end{itemize}

\begin{center}
\begin{tabular}{|@{}c@{}||@{}c@{}||@{}c@{}||@{}c@{}||@{}c@{}||@{}c@{}||@{}c@{}||@{}c@{}|}
\hline
1 & 2 &  3 & 4 & 5 & 6 & 7 & 8 \\
\hline
\begin{tabular}{cccc}
\\\textcircled{א}\\\\
\textcircled{ב}\\\\
\textcircled{ג}\\\\
\textcircled{ד}\\\\
\end{tabular}
&
\begin{tabular}{cccc}
\\\textcircled{א}\\\\
\textcircled{ב}\\\\
\textcircled{ג}\\\\
\textcircled{ד}\\\\
\end{tabular}

&
\begin{tabular}{cccc}
\\\textcircled{א}\\\\
\textcircled{ב}\\\\
\textcircled{ג}\\\\
\textcircled{ד}\\\\
\end{tabular}
&
\begin{tabular}{cccc}
\\\textcircled{א}\\\\
\textcircled{ב}\\\\
\textcircled{ג}\\\\
\textcircled{ד}\\\\
\end{tabular}
&
\begin{tabular}{cccc}
\\\textcircled{א}\\\\
\textcircled{ב}\\\\
\textcircled{ג}\\\\
\textcircled{ד}\\\\
\end{tabular}
&
\begin{tabular}{cccc}
\\\textcircled{א}\\\\
\textcircled{ב}\\\\
\textcircled{ג}\\\\
\textcircled{ד}\\\\
\end{tabular}
&
\begin{tabular}{cccc}
\\\textcircled{א}\\\\
\textcircled{ב}\\\\
\textcircled{ג}\\\\
\textcircled{ד}\\\\
\end{tabular}
&
\begin{tabular}{cccc}
\\\textcircled{א}\\\\
\textcircled{ב}\\\\
\textcircled{ג}\\\\
\textcircled{ד}\\\\
\end{tabular}
\\
\hline 
\begin{tikzpicture}
  \draw[dotted] (0,0) circle (0.8cm);
\end{tikzpicture}
& 
 \begin{tikzpicture}
  \draw[dotted] (0,0) circle (0.8cm);
\end{tikzpicture}
 &
 \begin{tikzpicture}
  \draw[dotted] (0,0) circle (0.8cm);
\end{tikzpicture}
 &
 \begin{tikzpicture}
  \draw[dotted] (0,0) circle (0.8cm);
\end{tikzpicture}
 &
 \begin{tikzpicture}
  \draw[dotted] (0,0) circle (0.8cm);
\end{tikzpicture}
 &
 \begin{tikzpicture}
  \draw[dotted] (0,0) circle (0.8cm);
\end{tikzpicture}
 &
 \begin{tikzpicture}
  \draw[dotted] (0,0) circle (0.8cm);
\end{tikzpicture}
 &
 \begin{tikzpicture}
  \draw[dotted] (0,0) circle (0.8cm);
\end{tikzpicture}
   \\
 \hline
\end{tabular}
\end{center}


\subsection*{חלק ב': שאלות פתוחות}
\begin{itemize}
    \item יש לבחור בדיוק 3 מתוך 4 שאלות פתוחות. משקל כל שאלה 20 נקודות.

    \item יש לסמן את השאלות שנבחרו בטבלה בדף זה. שאלות שלא סומנו לא תבדקנה!
    \item יש לנמק את התשובות בחלק ב' ולכתוב את כל החישובים הרלוונטים.
    \item יש לענות על כל השאלות בחלק זה על טופס המבחן במקום המיועד אחרי כל שאלה.   
    
    \underline{
    ניתן גם לכתוב מצדו השני של הדף. 
    }
    מחברת הטיוטה לא תבדק!
\end{itemize}
\bigskip

 בחלק ב' יש לבדוק את השאלות הבאות. סמנו את התשובות המתאימות --  3 שאלות לכל היותר!


\begin{center}
\setlength{\arrayrulewidth}{0.6pt}
\renewcommand{\arraystretch}{1.4}
\begin{tabular}{|c|c|c|c|}
\hline
\textcircled{1} & \textcircled{2} & \textcircled{3} & \textcircled{4}\\
\hline
\end{tabular}
\end{center}

\begin{center}
\huge{\textbf{בהצלחה!}}
\end{center}
\pagebreak
\subsection*{הוראות למילוי הטבלה}
\begin{itemize}
    \item[$\bullet$] יש להשתמש בעט שחור או כחול. 
    \item[$\bullet$] יש לסמן תשובה אחת בלבד ע"י השחרת העיגול המתאים.
    \item[$\bullet$]
    בשאלה בה לא יושחר אף עיגול,
    הניקוד יהיה 0.
    \item[$\bullet$] בשאלה בה יושחרו יותר מעיגול אחד, הניקוד יהיה 0.
.
    \item[$\bullet$] כדי לתקן סימון, יש לסמן X
    על התשובה הלא נכונה ולהשחיר את התשובה הנכונה.

    \item[$\bullet$]
    \textbf{שימו לב:}
    לא ניתן לבחור סעיף שסומן עליו X
    לכן הקפידו לסמן את תשובתכם כשאתם בטוחים בה.
\end{itemize}

\subsection*{דוגמאות}


\begin{center}
\begin{tabular}{|c||c||c||c||c||c||c||c||c||c|}
\hline
1 & 2 &  3 & 4 & 5 & 6 & 7 & 8 & 9 & 10 \\
\hline
\begin{tabular}{c}
\\\Large{\CIRCLE}\\\\
\ \  \hspace{-8pt}\raisebox{0pt}[0pt][0pt]{\Large\CIRCLE}%
\  \hspace{4pt}\llap{\raisebox{-4pt}[0pt][0pt]{ $\scalebox{3}{$\times$}$}}\hspace{-8pt}
\\\\
\textcircled{ג}\\\\
\textcircled{ד}\\\\
\end{tabular}
&
\begin{tabular}{cccc}
\\\textcircled{א}\\\\
\Large{\CIRCLE}\\\\
\textcircled{ג}\\\\
\textcircled{ד}\\\\
\end{tabular}
&
\begin{tabular}{cccc}
\\
\textcircled{א}
\\\\
\ \  \hspace{-8pt}\raisebox{0pt}[0pt][0pt]{\Large\CIRCLE}%
 \ \hspace{4pt}\llap{\raisebox{-4pt}[0pt][0pt]{ $\scalebox{3}{$\times$}$}}\hspace{-8pt}
\\\\
\Large{\CIRCLE}\\\\
\ \  \hspace{-8pt}\raisebox{0pt}[0pt][0pt]{\Large\CIRCLE}%
 \ \hspace{4pt}\llap{\raisebox{-4pt}[0pt][0pt]{ $\scalebox{3}{$\times$}$}}\hspace{-8pt}
\\\\
\end{tabular}
&
\begin{tabular}{cccc}
\\\textcircled{א}\\\\
\textcircled{ב}\\\\
\textcircled{ג}\\\\
\textcircled{ד}\\\\
\end{tabular}
&
\begin{tabular}{cccc}
\\\hspace{-8pt}\ \  \raisebox{0pt}[0pt][0pt]{\Large\CIRCLE}%
 \ \hspace{4pt}\llap{\raisebox{-4pt}[0pt][0pt]{ $\scalebox{3}{$\times$}$}}\hspace{-8pt}
\\\\
\textcircled{ב}\\\\
\textcircled{ג}\\\\
\textcircled{ד}\\\\
\end{tabular}
&
\begin{tabular}{cccc}
\\\hspace{-8pt}\ \  \raisebox{0pt}[0pt][0pt]{\Large\CIRCLE}%
 \ \hspace{4pt}\llap{\raisebox{-4pt}[0pt][0pt]{ $\scalebox{3}{$\times$}$}}\hspace{-8pt}
\\\\
\hspace{-8pt}\ \  \raisebox{0pt}[0pt][0pt]{\Large\CIRCLE}%
 \ \hspace{4pt}\llap{\raisebox{-4pt}[0pt][0pt]{ $\scalebox{3}{$\times$}$}}\hspace{-8pt}
\\\\
\textcircled{ג}\\\\
\textcircled{ד}\\\\
\end{tabular}
&
\begin{tabular}{cccc}
\\\Large{\CIRCLE}\\\\
\Large{\CIRCLE}\\\\
\textcircled{ג}\\\\
\textcircled{ד}\\\\
\end{tabular}
&
\begin{tabular}{cccc}
\\\Large{\CIRCLE}\\\\
\Large{\CIRCLE}\\\\
\Large{\CIRCLE}\\\\
\textcircled{ד}\\\\
\end{tabular}
&
\begin{tabular}{cccc}
\\\hspace{-8pt}\ \  \raisebox{0pt}[0pt][0pt]{\Large\CIRCLE}%
 \ \hspace{4pt}\llap{\raisebox{-4pt}[0pt][0pt]{ $\scalebox{3}{$\times$}$}}\hspace{-8pt}
\\\\
\Large{\CIRCLE}\\\\
\textcircled{ג}\\\\
\Large{\CIRCLE}\\\\
\end{tabular}
&
\begin{tabular}{cccc}
\\\hspace{-8pt}\ \  \raisebox{0pt}[0pt][0pt]{\Large\CIRCLE}%
 \ \hspace{4pt}\llap{\raisebox{-4pt}[0pt][0pt]{ $\scalebox{3}{$\times$}$}}\hspace{-8pt}
\\\\
\Large{\CIRCLE}
\\\\
\hspace{-8pt}\ \  \raisebox{0pt}[0pt][0pt]{\Large\CIRCLE}%
 \ \hspace{4pt}\llap{\raisebox{-4pt}[0pt][0pt]{ $\scalebox{3}{$\times$}$}}\hspace{-8pt}
\\\\
\Large{\CIRCLE}\\\\
\end{tabular}
\\
\hline
\end{tabular}
\ 

\ 

\begin{itemize}[label=$\bullet$]
    \item שאלות 1,2,3 עם סימונים תקינים עבור תשובות:
    \begin{enumerate}
        \item א
        \item ב
        \item ג
    \end{enumerate}
    \item בשאלות 4,5,6 חסר סימון והניקוד יהיה 0.
    \item בשאלות 7,8,9,10
    יש יותר מסימון אחד והניקוד יהיה 0
\end{itemize}


\end{center}


%%%
%%%
%%%   End of Header and Instructions
%%%
%%%

\pagebreak
%%%%%%%%%%%%%%%%%%%%%%%%%%%%%%%%%%%%%%
%%  Problems
%%%%%%%%%%%%%%%%%%%%%%%%%%%%%%%%%%%%%%


%%%%% שאלה 2 %%%%%%%

\section*{
חלק א' 
(40 נק')
להשחיר את התשובות בטבלה!}
\begin{enumerate}
\item\points{5}
יהיו
$H_1,H_2$
מישורים הנתונים על ידי 
$$,H_1=\{
(x,y,z)\in\rr^3 \bigm| x-z=2\},\quad H_2=\{
(1+t, 1-t+s, s)
\bigm| t,s\in\rr\}$$
עבור
הישר
$\ell= H_1\cap H_2$.
מה מהסעיפים הבאים \underline{\textbf{לא}} נכון.
\begin{enumerate}
\item
קיים
$x$
עבורו
$(x,0,0)\in \ell$
\item
קיים
$y$
עבורו
$(0,y,0)\in \ell$
\item
קיים
$z$
עבורו
$(0,0,z)\in \ell$
\item
קיים
$t$
עבורו
$(t,0,-t)\in \ell$
\end{enumerate}

\noindent\rule{\textwidth-30pt}{0.4pt}


\item\points{5}
עבור
$a\in\rr$
נתון, נסמן
$$A=\begin{pmatrix}
a&1&0\\
0&a&1\\
1&0&a\\
\end{pmatrix}$$
בחרו את הטענה הנכונה:
\begin{enumerate}
 %%%1
\item
עבור
$a=-1$
המטריצה
$A$
הפיכה.
%%%
\item
עבור
$a=0$
מתקיים
$\mathrm{rank}(A)=2$.
%%%
\item
עבור
$a=2$
מתקיים
$\dim(N(A))=\dim(Row(A))$.
%%%%%%
\item
עבור
$a=-3$
מתקיים
$\det(A)\neq0$.
%%%
\end{enumerate}

\noindent\rule{\textwidth-30pt}{0.4pt}

\item\points{5}
תהי $T:\rr^2\to \rr^2$ העתקה לינארית המקיימת
$T(1,1) = (1,2)$.
איזה סעיף בהכרח נכון?
\begin{enumerate}
\item $\im(T)=\ker(T)$
\item $\dim(\im(T))\geq \dim(\ker(T))$
\item $\im(T)\neq \ker(T)$
\item $\dim(\im(T))\leq \dim(\ker(T))$
\end{enumerate}
\noindent\rule{\textwidth-30pt}{0.4pt}


%%%%% שאלה 2 %%%%%%%

\item\points{5}
יהיו
$a,b\in\rr$
כך שהקבוצה
$$\left\{
\begin{pmatrix}
a\\1\\1\\
\end{pmatrix},
\begin{pmatrix}
2\\-1\\1-2a\\
\end{pmatrix},
\begin{pmatrix}
2\\b\\-2\\
\end{pmatrix}
\right\}$$
היא בסיס אורתוגונלי של
$\rr^3$.
למה שווה
$a+b$\ ?
\begin{enumerate}
\begin{multicols}{4}
%%%3
\item
$2$
\smallskip
%%%2
\item
$0$
\smallskip
%%%4
\item
$4$
\smallskip
%%%1
\item
$3$
\smallskip
\end{multicols}
\end{enumerate}
%\noindent\rule{\textwidth-30pt}{0.4pt}
\pagebreak

\item\points{5}
נתון 
$B=\{v_1,v_2\}$
בסיס סדור של
$\rr^2$
עבורו מתקיים
\begin{align*}
\left[\begin{pmatrix}
1\\0\\
\end{pmatrix}\right]_B=
\begin{pmatrix}
1\\3 \\
\end{pmatrix}\quad,\quad
\left[\begin{pmatrix}
0\\1\\
\end{pmatrix}\right]_B=
\begin{pmatrix}
2\\5 \\
\end{pmatrix}
\end{align*}
למה שווה
$v_1+v_2$:
\begin{enumerate}
\begin{multicols}{4}
%%%1
\item
$
\begin{pmatrix}
1\\4\\
\end{pmatrix}$
%%%3
\item
$
\begin{pmatrix}
0\\-2\\
\end{pmatrix}$
%%%4
\item
$
\begin{pmatrix}
-3\\
2\\
\end{pmatrix}$
%%%2
\item
$
\begin{pmatrix}
-1\\1\\
\end{pmatrix}$\end{multicols}
\end{enumerate}

\noindent\rule{16cm}{0.4pt}
%\pagebreak

%%%%
\item\points{5}
נתון כי
$A,B\in M_{\ntn}(\rr)$
מטריצות שקולות שורה.
איזו מהטענות הבאות 
\underline{ בהכרח נכונה:}
\begin{enumerate}
%%%2
\item
$\det(A)=\det(B)$
%%%1
\item
$N(A)=N(B)$
%%%4
\item
$Av=Bv$
לכל
$v\in\rr^n$.
%%%3
\item
$A^2,B^2$
שקולות שורה.
\end{enumerate}



\noindent\rule{16cm}{0.4pt}
\item\points{5}
נתון כי
$A,B\in M_{3\times 3}(\rr)$
מטריצות מדורגות קנונית.
איזו מהטענות הבאות 
\underline{בהכרח נכונה:}
\begin{enumerate}
%%%4
\item
$\det(A)+\det(B)\le 2$
%%%2
\item
$A+B$
מדורגת קנונית.
%%%1
\item
$\det(A)\cdot \det(B)=0$
%%%3
\item
למטריצה
$A-B$
יש לפחות שורת אפסים אחת.
\end{enumerate}

\noindent\rule{16cm}{0.4pt}


\item\points{5}
תהי 
$A\in M_{4\times 4}(\rr)$
כך ש-
$\rank(A)=1$. איזו מהטענות הבאות 
\underline{בהכרח נכונה}?
\begin{enumerate}
    \item למטריצה $A$
    יש 
    לפחות
    2 ערכים עצמיים שונים.
    \item המטריצה $A$ לכסינה.
        \item המטריצה $A$ אינה לכסינה.
         \item למטריצה $A$ יש לכל היותר 2 ערכים עצמיים שונים.
\end{enumerate}


\pagebreak

\end{enumerate}
%%%
%%%%%%
%%%

 \pagebreak

\section*{חלק ב' (60 נק') - יש לבחור 3 מתוך 4 ולסמן בעמוד הראשון.}
יש לענות במקומות המתאימים בעמודים הבאים.
\begin{enumerate}
\newcommand{\firstQ}{
\item\points{20}
יהי
$a\in\rr$
ותהי המטריצה
$$.A=
\begin{pmatrix}
a&2&-1\\
0&a&0\\
-a&-2&1\\
\end{pmatrix}\in
M_{3\times3}(\rr)$$
\begin{enumerate}
%%%%%%% חלק ב, שאלה 1 לכסון, סעיף א %%%%%%%
\item\points{4}
מצאו את הפולינום האופייני של
$A$.
%%%%%%%חלק ב, שאלה 1 לכסון, סעיף ב %%%%%%%
\item
\points{12}
לכל ערך של
$a\in\rr$
קבעו האם 
$A$
לכסינה, והסבירו את נימוקיכם/ן.
%%%%%%חלק ב, שאלה 1 לכסון, סעיף ג %%%%%%
\item\points{4}
עבור
$a=0$
חשבו את
$A^{5777}$.
\end{enumerate}

}
\firstQ
\noindent\rule{\textwidth-30pt}{0.4pt}

\newcommand{\secondQ}{
\item\points{20}
נסמן את הבסיס 
$B=\{v_1,v_2,v_3\}$
של
$\rr^3$,
הנתון על ידי הוקטורים
$$v_1=\begin{pmatrix}
2\\1\\1\\
\end{pmatrix},\,
v_2=\begin{pmatrix}
1\\1\\1\\
\end{pmatrix},\,
v_3=\begin{pmatrix}
1\\0\\-1\\
\end{pmatrix}$$
נתונה העתקה לינארית
$T:\rr^3\to\rr^3$
המקיימת
 $$
 .T(v_1)=
 \begin{pmatrix}
1\\1\\0\\
 \end{pmatrix},\,
 T(v_2)=
  \begin{pmatrix}
1\\-1\\2\\
 \end{pmatrix},
\,
  T(v_3)=
 \begin{pmatrix}
3\\1\\2\\
 \end{pmatrix}
$$
\begin{enumerate}
\item\points{8}
מצאו נוסחה עבור
$T\begin{pmatrix}
a\\b\\c\\
\end{pmatrix}$.
%%%
\item\points{8}
מצאו בסיס אורתונורמלי לגרעין 
$\ker(T)$
ומצאו בסיס אורתונורמלי לתמונה
$\im(T)$.
\item\points{4}
האם קיים בסיס
סדור
$C$
של
$\rr^3$
עבורו מתקיים
$$?[T]_C^C=
\begin{pmatrix}
0&0&1\\
0&0&2\\
0&0&3\\
\end{pmatrix}$$
%%%%%
%%%%%%חלק ב, שאלה 2 העתקות לינאריות, סעיף ג %%%%%%
\end{enumerate}
}
\secondQ

%\noindent\rule{\textwidth-30pt}{0.4pt}
\pagebreak

\newcommand{\thirdQ}{
\item\points{20}
תהי
$$.A=\begin{pmatrix}
1&2&-1\\
-1&1&1\\
-2&8&2\\
\end{pmatrix}
\in M_{3\times 3}(\rr)$$
יהיו
$U,W$
תתי-המרחבים של
$\rr^3$
המוגדרים על ידי
$$.W=\Set{A\cdot v | v\in \rr^3},\quad
U=N(A)
$$
\begin{enumerate}
%%%%%%% חלק ב, שאלה 3 ארותוגונליות, סעיף א %%%%%%%
\item\points{6}
מצאו בסיס ומימד של
$W$.
%%%
\item\points{5}
מצאו בסיס ומימד של
$U$.
%%%
%%%
\item\points{6}
מצאו בסיס ומימד של
$U\cap W$ ושל $U+W$.
%%%%%
\item
\points{3}
האם קיים וקטור
$0\neq v\in\rr^3$
המקיים
$v\in Row(A)\cap Col(A)$?
%%%%%%%חלק ב, שאלה 3 ארותוגונליות, סעיף ב %%%%%%%
\end{enumerate}

}
\thirdQ

\noindent\rule{\textwidth-30pt}{0.4pt}
\newcommand{\forthQ}{
    \item \points{20}
    הסעיפים הבאים לא קשורים אחד לשני:
    \begin{enumerate}
        \item\points{10} 
        תהי $T:V\to W$ \ 
        העתקה לינארית
        ויהיו
        $v_1,v_2,v_3\in V$.
        
        הוכיחו כי אם
        $T(v_1),T(v_2),T(v_3)$
        בלתי תלויים לינארית (בת"ל),
        
        אז גם
        $v_1,v_2,v_3$
        בלתי תלויים לינארית (בת"ל).

        \ 
        
        \item\points{10} 
        תהי מטריצה $A\in M_{n\times n}(\cc)$ כך ש- $A^2=A$ ויהי
        $\lambda$
        ערך עצמי של $A$.

        הראו כי
        $\lambda=1$ או $\lambda = 0$.
    
    \end{enumerate}
}\forthQ
\pagebreak
\setcounter{enumi}{0}
    \firstQ
\noindent\rule{\textwidth-30pt}{0.4pt}
\underline{פיתרון:}
    
    \pagebreak \ 
    \ 
    
    \pagebreak

    
    \secondQ
    \noindent\rule{\textwidth-30pt}{0.4pt}
    \underline{פיתרון:}


    \pagebreak \ \pagebreak \ 

    \ 
    
    \pagebreak

    \thirdQ
    \noindent\rule{\textwidth-30pt}{0.4pt}
    \underline{פיתרון:}


    \pagebreak \ \pagebreak

    \forthQ
    \noindent\rule{\textwidth-30pt}{0.4pt}
    \underline{פיתרון:}


    \pagebreak \ \pagebreak
\end{enumerate}

    \pagebreak 
    
    עמוד נוסף לכתיבת המשך פיתרונות
    
    \pagebreak
    עמוד נוסף לכתיבת המשך פיתרונות

 \pagebreak

 % --- דף הנוסחאות ---

\section*{דף הנוסחאות}

\begin{enumerate}
\item $M_{m\times n}(F)$ זו קבוצת המטריצות בגודל $n\times m$
($m$ שורות, $n$ עמודות)
עם רכיבים מ-$F$.

\noindent\rule{\textwidth-30pt}{0.4pt}

\item לכל מטריצה $A$, נסמן על ידי $(A)_{i,j}$ את האיבר של $A$ הנמצא בשורה ה-$i$ ובעמודה ה-$j$.

\noindent\rule{\textwidth-30pt}{0.4pt}

\item פיתוח דטרמיננטה של $A\in M_{n\times n}(F)$ לפי השורה ה-$i$ או לפי העמודה ה-$j$:
\[
\det(A)=\sum_{i=1}^{n}(-1)^{i+j}a_{ij}\det\!\big(M_{ij}(A)\big)
=\sum_{j=1}^{n}(-1)^{i+j}a_{ij}\det\!\big(M_{ij}(A)\big)
\]
כאן $M_{ij}(A)$ זה המינור ה-$i,j$ של המטריצה $A$.

\noindent\rule{\textwidth-30pt}{0.4pt}

\item אם $T:V\to W$ העתקה לינארית, אם $B=\{v_1,\ldots,v_n\}$ בסיס סדור של $V$ ואם $C$ בסיס סדור של $W$,
אז המטריצה המייצגת של $T$ ביחס לבסיסים $B$ ו-$C$ מוגדרת על ידי
\[
[T]^B_C=
\begin{pmatrix}
\vline & & \vline\\
[T(v_1)]_C & \cdots & [T(v_n)]_C\\
\vline & & \vline
\end{pmatrix}.
\]

\noindent\rule{\textwidth-30pt}{0.4pt}

\item לכל $u=(a_1,\ldots,a_n),\; w=(b_1,\ldots,b_n)\in\mathbb{R}^n$

נסמן את המכפלה הפנימית הסטנדרטית ב-$\mathbb{R}^n$ (המכפלה הסקלרית):
\[
 u\cdot w=a_1\cdot b_1+\cdots+a_n\cdot b_n.
\]

\noindent\rule{\textwidth-30pt}{0.4pt}

% --- דף הנוסחאות ---


\item תהליך גראם-שמיט (גרסה א):
אם $W$ תת-מרחב של $\rr^n$, ואם $\{v_1,\ldots,v_m\}$ בסיס של $W$

אז $\{w_1,\ldots,w_m\}$ בסיס אורתוגונלי של $W$ כאשר
\[
w_1=v_1,\qquad
w_k=v_k-\sum_{i=1}^{k-1}\frac{ v_k\cdot w_i}{ w_i\cdot w_i}\,w_i
\quad (k=2,\ldots,m).
\]
\item תהליך גראם-שמיט (גרסה ב):
בהנתן סדרה בת"ל
$\{v_1,\ldots,v_k\}\subseteq \rr^n$, נבנה סדרה אורתונורמלית באופן הבא:
\[
w_1=\frac{v_1}{\|v_1\|},\qquad
w_i'=v_i-\sum_{l=1}^{i-1} (v_i\cdot w_l)\,w_l,\qquad w_i = \frac{w_i'}{\|w_i'\|}
\]

\noindent\rule{\textwidth-30pt}{0.4pt}

\item השורשים של המשוואה הריבועית
$
0=c+bx+ax^2,\qquad a\ne 0,
$
נתונים על ידי
\[
x_{1,2}=\frac{-b\pm\sqrt{b^2-4ac}}{2a}.
\]

\noindent\rule{\textwidth-30pt}{0.4pt}

\item מטריצה $A\in M_{n\times n}(\ff)$ דומה למטריצה $B\in M_{n\times n}(\ff)$ 

אם
קיימת 
מטריצה הפיכה
$P\in M_{n\times n}(\ff)$ כך ש-
$A=PBP^{-1}$.
\end{enumerate}

\end{document}